\section{Cyclic Artin--Schreier extensions of the Laurent field}

Throughout this section, $p$ is prime, $k=\mathbf F_p$, and
$K=k((t))$. Write $v$ for the additive order valuation, normalized by
$v(t)=1$ and $v(0)=+\infty$, and write $\wp(x)=x^p-x$.
The following construction supplies the cyclic case of the
Artin--Schreier correspondence used in \cite{Camina1997}, Theorem~1,
and makes explicit the ramification assertion for reduced parameters
in her Lemma~3. The valuation is constructed on a polynomial basis;
existence of an extension of $v$ is not assumed.

\begin{lemma}[A pole prime to the characteristic is not an Artin--Schreier image]
\label{lem:artin-schreier-no-root-from-pole}
\source{camina1997:page-0005}
\notready
If $\gamma\in K$, $m>0$, $p\nmid m$, and $v(\gamma)=-m$, then
$\gamma\notin\wp(K)$.
\end{lemma}

\begin{proof}
Suppose $\gamma=x^p-x$. If $v(x)\ge0$, then
$v(x^p-x)\ge\min\{p v(x),v(x)\}\ge0$, a contradiction.
Thus $v(x)<0$. Since $p\ge2$, we have $p v(x)<v(x)$, so the
two summands have different valuations and
$v(x^p-x)=p v(x)$. This gives $-m=p v(x)$, contradicting
$p\nmid m$.
\end{proof}

\begin{lemma}[Irreducibility of an Artin--Schreier polynomial]
\label{lem:artin-schreier-irreducible}
\source{camina1997:page-0003}
\notready
Let $F$ be any field of characteristic $p$ and $a\in F$.
If $a\notin\wp(F)$, then $Y^p-Y-a$ is irreducible in $F[Y]$.
\end{lemma}

\begin{proof}
Choose a root $y$ in an algebraic closure of $F$. The $p$ distinct
elements $y+c$, for $c\in\mathbf F_p$, are roots, because
$(y+c)^p-(y+c)=y^p-y$. They are all the roots, since the
polynomial has degree $p$; they are simple, since its derivative
is $-1$.
Let $q$ be the monic minimal polynomial of $y$ and $d=\deg q$.
It divides $Y^p-Y-a$, so $1\le d\le p$ and it has distinct roots
$y+c_1,\ldots,y+c_d$ for some $c_i\in\mathbf F_p$.
Its coefficient of $Y^{d-1}$ is therefore
$-d y-\sum_i c_i$. This coefficient and $\sum_i c_i$ belong
to $F$. If $d<p$, the element $d\cdot1$ is nonzero in $F$,
so this coefficient identity gives $y\in F$, contrary to
$a\notin\wp(F)$. Hence $d=p$. The two monic polynomials
$q$ and $Y^p-Y-a$ have equal degree and one divides the other;
they are equal. This proves irreducibility.
\end{proof}

\begin{lemma}[The cyclic extension defined by a reduced pole]
\label{lem:artin-schreier-cyclic-extension}
\uses{lem:artin-schreier-no-root-from-pole,lem:artin-schreier-irreducible}
\source{camina1997:page-0003,camina1997:page-0005}
\notready
Suppose $v(\gamma)=-m<0$ and $p\nmid m$. The quotient
\[
 L=K[Y]/(Y^p-Y-\gamma)
\]
is a field of degree $p$ over $K$. Writing $y$ for the class of
$Y$, every element of $L$ has a unique expression
$\sum_{i=0}^{p-1}a_i y^i$ with $a_i\in K$.
The maps $\sigma_c(y)=y+c$, fixing $K$, identify
$(\mathbf F_p,+)$ with $\operatorname{Gal}(L/K)$.
\end{lemma}

\begin{proof}
The first two lemmas give irreducibility. Division by the monic
polynomial gives a unique remainder of degree less than $p$,
and hence the displayed basis and degree. The quotient is a
field: for a nonzero remainder $h$, irreducibility gives
$\gcd(h,Y^p-Y-\gamma)=1$, and a Bezout identity gives an inverse
to its class.
Substitution $y\mapsto y+c$ respects the defining relation,
because $c^p=c$, and substitution by $-c$ is its inverse.
These automorphisms satisfy
$\sigma_c\sigma_d=\sigma_{c+d}$ and are distinct by their
values on $y$. Every $K$-automorphism sends $y$ to a root of
$Y^p-Y-\gamma$, whose roots are exactly these $y+c$.
Thus they are all the automorphisms. The polynomial splits in
$L$ and has derivative $-1$, so the extension is normal and
separable, as required for the stated Galois group.
\end{proof}

\begin{lemma}[The valuation on the cyclic extension]
\label{lem:artin-schreier-weighted-valuation}
\uses{lem:artin-schreier-cyclic-extension}
\source{camina1997:page-0005,camina1997:page-0006}
\notready
For $L$, $y$, $m$ as above, the formula
\[
 w\left(\sum_{i=0}^{p-1}a_i y^i\right)
   =\min_{a_i\ne0}\{p v(a_i)-mi\},\qquad w(0)=+\infty,
\]
defines a surjective additive discrete valuation on $L$.
It satisfies $w|_K=p v$ and $w(y)=-m$.
\end{lemma}

\begin{proof}
The formula is well-defined by uniqueness of the polynomial
expression. Distinct nonzero terms have distinct weights: an
equality $p v(a_i)-mi=p v(a_j)-mj$ implies $p\mid m(i-j)$,
hence $p\mid i-j$ and then $i=j$, since $0\le i,j<p$.
In particular every nonzero expression has a unique least-weight
term. The inequalities for $v$ on each coefficient give
$w(x+x')\ge\min\{w(x),w(x')\}$, and clearly $w(1)=0$.

We prove multiplicativity. Let the least-weight terms of nonzero
$x=\sum a_i y^i$ and $z=\sum b_j y^j$ have indices $i_0,j_0$.
Among the nonzero terms $a_i b_j y^{i+j}$ in the unreduced
product, exactly $(i_0,j_0)$ has least weight
$w(x)+w(z)$. If $i+j<p$, this term is already in basis form.
If $p\le i+j\le2p-2$, the relation $y^p=y+\gamma$ gives
\[
 a_i b_j y^{i+j}
   =a_i b_j y^{i+j-p+1}
      +a_i b_j\gamma y^{i+j-p}.
\]
Both exponents on the right lie between $0$ and $p-1$.
The second term has the original weight
$p v(a_i)+p v(b_j)-m(i+j)$, because $v(\gamma)=-m$;
the first has that weight plus $m(p-1)>0$.
Thus reduction leaves exactly one contribution of globally least
weight, namely the lower-weight contribution of $(i_0,j_0)$.
Contributions to the same basis coefficient with larger weight
have larger coefficient valuation, so they cannot cancel its
leading coefficient. Consequently
$w(xz)=w(x)+w(z)$. The assertion also holds when a factor is zero.
Applying it to $xx^{-1}=1$ gives $w(x^{-1})=-w(x)$.

The formula gives $w(t)=p$ and $w(y)=-m$. Since $p\nmid m$,
there are integers $a,b$ with $ap-bm=1$; the nonzero element
$t^a y^b$ has valuation one. Its integral powers realize every
integer, proving surjectivity. The formula on constant basis
expressions gives $w|_K=p v$.
\end{proof}

\begin{lemma}[Completeness and residue field]
\label{lem:artin-schreier-complete-residue}
\uses{lem:artin-schreier-weighted-valuation}
\source{camina1997:page-0005,camina1997:page-0006}
\notready
The valued field $(L,w)$ is complete. Its residue field is
canonically $k$, and the ramification index of $L/K$ is $p$.
In particular $L/K$ is totally ramified.
\end{lemma}

\begin{proof}
The formula for $w$ shows that a sequence
$x_n=\sum_{i<p}a_{n,i}y^i$ is Cauchy precisely when each of its
$p$ coefficient sequences is Cauchy in $K$: a lower bound on
$w(x_n-x_r)$ is exactly the corresponding lower bound on all
$p v(a_{n,i}-a_{r,i})-mi$.
Each coefficient sequence converges, since $k((t))$ is complete.
Indeed, for a Cauchy sequence of Laurent series the coefficients
in each bounded range eventually stabilize; their stabilized
values have a common lower support bound, obtained by comparing
the tail of the sequence with one fixed term. They therefore
define a Laurent series, and the Cauchy condition gives
convergence to it. Taking these limits for all $i<p$ and applying
the formula for $w$ proves convergence of $x_n$ in $L$.
The valuations define metrics by $d(x,z)=2^{-w(x-z)}$ and
$2^{-v(x-z)}$, with distance zero for equal arguments, so this
sequential completeness is completeness of the valuation
uniformity.

Let $x=\sum_{i<p}a_i y^i$ satisfy $w(x)\ge0$.
Each nonzero basis term has nonnegative weight. For $i>0$
that weight cannot be zero, since $p\nmid mi$; all such terms
lie in the maximal ideal $\{z:w(z)>0\}$. The constant basis
coefficient $a_0$ has $v(a_0)\ge0$. If $c\in k$ is its
constant coefficient as a power series in $t$, then
$v(a_0-c)>0$, so $w(x-c)>0$.
The class of every valuation-ring element is therefore represented
by a constant in $k$. A nonzero constant has valuation zero,
so distinct constants give distinct residue classes. Reduction
on constants is a field homomorphism, proving the claimed
residue-field identification. Finally
$w(L^\times)=\mathbf Z$ and $w(K^\times)=p\mathbf Z$;
their index is $p=[L:K]$. Together with the trivial residue
extension, this is total ramification.
\end{proof}

\begin{lemma}[A Laurent-series parameter for the extension]
\label{lem:artin-schreier-laurent-parameter}
\uses{lem:artin-schreier-complete-residue}
\source{camina1997:page-0005}
\notready
For every $\pi\in L$ with $w(\pi)=1$, the map
\[
 k((T))\longrightarrow L,\qquad
 \sum_{n\ge n_0}c_n T^n\longmapsto\sum_{n\ge n_0}c_n\pi^n
\]
is an isomorphism of topological $k$-algebras. It identifies
the order in $T$ with $w$.
\end{lemma}

\begin{proof}
The sum converges by completeness, since
$w(c_n\pi^n)\ge n$ and the negative-exponent part is finite.
If the first nonzero coefficient has index $r$, the term
$c_r\pi^r$ has valuation $r$ and the remaining tail has
valuation at least $r+1$. The latter bound holds for the limit:
if a convergent sequence has valuations at least $r+1$,
then sufficiently small differences from its limit give the
same bound on the limit. The valuation of the sum is thus $r$.
This proves injectivity and the asserted preservation of order.
Finite sums preserve addition and multiplication. Their tails
tend to zero; for products of two Laurent series the omitted
terms also tend to zero, since both series have lower support
bounds. Passing to the limits proves that the map preserves
addition and multiplication, and it plainly fixes $k$.

To prove surjectivity, first take $x$ with $w(x)\ge0$.
Put $x_0=x$. The residue-field assertion supplies a unique
$c_n\in k$ with $w(x_n-c_n)>0$; set
$x_{n+1}=(x_n-c_n)/\pi$. Then $w(x_{n+1})\ge0$ and induction gives
\[
 x=\sum_{j=0}^{N-1}c_j\pi^j+\pi^N x_N.
\]
The remainder has valuation at least $N$, so the partial sums
converge to $x$. An arbitrary element of $L$ becomes integral
after multiplication by a sufficiently large power of $\pi$,
giving its Laurent expansion as well. Finally preservation of
order makes both this map and its inverse continuous for the
valuation topologies.
\end{proof}

\begin{lemma}[Depth of an Artin--Schreier translation]
\label{lem:artin-schreier-translation-depth}
\uses{lem:artin-schreier-cyclic-extension,lem:artin-schreier-weighted-valuation,lem:artin-schreier-laurent-parameter}
\source{camina1997:page-0009}
\notready
Choose $1\le b<p$ and an integer $a$ with $ap-bm=1$, and put
$\pi=t^a y^b$. Then $w(\pi)=1$, every $\sigma_c$ is a
valuation-preserving continuous automorphism of $L$, and for
$c\in\mathbf F_p^\times$,
\[
 w(\sigma_c(\pi)-\pi)=m+1.
\]
\end{lemma}

\begin{proof}
Multiplication by $m$ permutes the nonzero residue classes modulo
$p$, so $b$ exists with $bm\equiv-1\pmod p$; taking
$a=(bm+1)/p$ gives the desired identity and $w(\pi)=1$.
For each basis term, expand
$\sigma_c(a_i y^i)=a_i(y+c)^i$.
Its degree-$i$ term is $a_i y^i$. Every other nonzero term has
degree $j<i$ and a nonzero constant binomial coefficient, so its
weight exceeds that of $a_i y^i$ by $m(i-j)>0$.
The unique least-weight term of an arbitrary nonzero basis
expression therefore remains unchanged by $\sigma_c$.
Thus $w(\sigma_c(x))=w(x)$ for every $x$, and the same equality
applied to differences proves continuity, including continuity
of the inverse $\sigma_{-c}$.

For $c\ne0$, put $u=c/y$, so $w(u)=m>0$. Since $\sigma_c$
fixes $t$,
\[
 \sigma_c(\pi)-\pi
   =\pi\bigl((1+u)^b-1\bigr).
\]
In the binomial expansion the linear term $bu$ has valuation
$m$, because $1\le b<p$. Every nonzero term of degree at least
two has valuation at least $2m>m$. The expression in parentheses
therefore has valuation $m$, proving the formula.
\end{proof}

\begin{theorem}[Cyclic subgroups with prescribed prime-to-$p$ depth]
\label{thm:nottingham-cyclic-prescribed-depth}
\uses{lem:artin-schreier-translation-depth,lem:nottingham-laurent-equivalence,def:nottingham-depth,lem:nottingham-depth-filtration}
\source{camina1997:page-0005,camina1997:page-0009}
\notready
For every positive $m$ with $p\nmid m$, the Nottingham group
$J_p$ contains a cyclic subgroup of order $p$ whose nonidentity
elements all have depth $m$.
\end{theorem}

\begin{proof}
Take $\gamma=t^{-m}$ and the field and parameter constructed
above. Transport the automorphisms $\sigma_c$ to $k((T))$
through the parameter isomorphism. They are continuous
$k$-algebra automorphisms, and their parameter differences have
order $m+1\ge2$ when $c\ne0$; for $c=0$ the difference is zero.
A Laurent series of order at least two is $T^2$ times a power
series. The Laurent--Nottingham equivalence therefore identifies
each transported automorphism with an element of $J_p$ and
respects their composition. The translations are distinct and
compose according to addition in $\mathbf F_p$, so their image
is cyclic of order $p$. The order formula for every nonzero $c$
gives exactly depth $m$.
\end{proof}
